%% ma: L11-B17-0024
%% lop: 11
%% chuong: 5
%% bai: 17
%% dang: 11.17.1
%% loai: TN4
%% muc: TH
%% dap_an: "A"
%% nguon: "tai-lieu-hoc-tap-mon-toan-lop-11-theo-chuong-trinh-giao-duc-pho-thong-2018.pdf (Lớp Toán Thầy Duy, Phạm Lê Duy, THPT Chu Văn An), Chương 3 Bài 3, Phần 1, Câu 4 (tr.399)"
%% kien_thuc: "Lớp 11 Bài 17 (hàm số liên tục)"
%% trang_thai: cho-duyet
%% ly_do: "Dạng mới đề xuất (11.17.x), chờ giáo viên duyệt dạng; tài liệu gốc không có lời giải mẫu, đáp án do kho tự giải và kiểm bằng sympy."
\begin{ex}
Cho hàm số $f(x)=\begin{cases}\dfrac{x-2}{\sqrt{x+2}-2}&\text{khi }x\ne2\\4&\text{khi }x=2\end{cases}$. Chọn mệnh đề đúng?
\choice{\True Hàm số liên tục tại $x=2$.}{Hàm số gián đoạn tại $x=2$.}{$f(4)=2$.}{$\lim\limits_{x\to2}f(x)=2$.}
\loigiai{\sloppy $\dfrac{x-2}{\sqrt{x+2}-2}=\sqrt{x+2}+2\to4=f(2)$ nên liên tục tại $x=2$ (A đúng, B và D sai). $f(4)=\dfrac2{\sqrt6-2}\ne2$.}
\end{ex}
